Krátký slovník klíčových pojmů určený pro pedagogy — rychlá reference při přípravě materiálů a komunikaci se studenty, s krátkými definicemi a praktickými poznámkami.

\begin{description}
  \item[Generativní AI] (general background knowledge) Algoritmy a modely, které vytvářejí nový obsah (text, obrázky, kód apod.) na základě naučených vzorů. V kontextu výuky jde především o nástroje, které studenti i vyučující mohou využít jako asistenty při tvorbě materiálů, vysvětlování či generování cvičení; je dobré rozlišovat tvorbu pod vedením učitele a plně automatické generování. Doporučuje se určit jasná pravidla použití a ověřování výsledků.
  \item[LLM — velké jazykové modely] (general background knowledge) Modely trénované na rozsáhlých korpusech textu, které předpovídají další tokeny a lze je použít pro generování textu, sumarizace či konverzační asistenci. V pedagogickém nasazení sledujte limity ve faktické přesnosti a vhodnosti obsahu. Učitelé by měli nastavit explicitní instrukce a validační kroky.
  \item[Tokeny] (general background knowledge) Základní jednotky textu, které model zpracovává (mohou být část slova, celé slovo nebo více slov v závislosti na tokenizéru). Důležité pro řízení délky vstupu/výstupu modelu a pro účtování spotřeby; při přípravě úloh plánujte rozpočet tokenů. V praxi to ovlivňuje délku promptů a rozsah poskytovaných kontextů.
  \item[Embeddings] (general background knowledge) Vektorové reprezentace textu/slov, které zachycují jejich význam a umožňují měřit podobnost mezi dokumenty; základní komponenta pro vyhledávání relevantních zdrojů nebo clustering odpovědí. Prakticky se používají pro vyhledávání v databázi učebních materiálů a pro návrh personalizovaných doporučení.
  \item[RAG — Retrieval‑Augmented Generation] (general background knowledge) Vzor, kdy se před generováním textu vyhledají relevantní dokumenty (retrieval) a tyto dokumenty se použijí jako kontext pro LLM, čímž se zvyšuje přesnost a faktická podpora odpovědí. (Viz také „Základní RAG blueprint pro kurzy“ v předchozí kapitole.) Vhodný při tvorbě fact‑checked materiálů a při omezování halucinací.
  \item[DPIA — posouzení dopadů na ochranu osobních údajů] (general background knowledge) Dokumentovaný postup hodnocení rizik zpracování osobních údajů (GDPR), doporučený před nasazením AI řešení, která pracují se studentskými daty. Obsahuje identifikaci rizik, návrh mitigací a záznamy o rozhodnutích. Výsledek DPIA pomáhá definovat technické a organizační opatření.
  \item[Minecraft Education] Výuková verze Minecraftu s připravenými světy a aktivitami vhodnými i pro ilustraci principů strojového učení a AI; některé světy pro „Hour of Code“ jsou zdarma, další jsou dostupné v rámci školních licencí Microsoft 365 A3/A5 \cite{kb_ai_101_tipu_pdf_04e4a115_page_15_2}. Lze jej využít pro projektové vyučování a interaktivní demonstrace, které zvyšují zapojení studentů.
  \item[Microsoft 365 Copilot / Copilot Chat] Chatbot založený na modelech GPT dostupný ve školních variantách Office 365; ve školním prostředí je k dispozici webová i mobilní verze Copilot Chat a placený doplněk Microsoft 365 Copilot rozšiřuje funkce přímo v aplikacích Word, Excel, PowerPoint apod. \cite{kb_ai_101_tipu_pdf_04e4a115_page_3_1}. Využití zahrnuje asistenční psaní, návrh prezentací a analýzu dat; při zavádění kontrolujte nastavení ochrany dat.
  \item["Pokrok v matematice"] Příklad nástroje pro generování a nasazení cvičných úloh, který umožňuje učiteli vybrat oblast, generovat příklady, nechat studenty odeslat postupy řešení; systém předopraví zadání a učitel finalizuje hodnocení, přičemž jsou dostupné přehledy výkonu jednotlivých žáků \cite{kb_ai_101_tipu_pdf_04e4a115_page_8_1}. Hodí se pro individualizovanou práci i zpětnou vazbu a pro monitoring pokroku třídy.
  \item[OneNote / math.microsoft.com] Praktické nástroje pro technické obory: OneNote umí převádět rukopisné rovnice do digitální podoby (včetně výpočtů a grafů) a web math.microsoft.com agreguje nástroje pro řešení matematických úloh \cite{kb_ai_101_tipu_pdf_04e4a115_page_8_1}. Tyto nástroje usnadňují sběr zadání a automatickou kontrolu postupů, což šetří čas učitele.
  \item[Halucinace (hallucination)] (general background knowledge) Situace, kdy model vygeneruje fakticky nepřesné nebo vymyšlené informace; riziko vyžaduje ověřování výstupů a lidskou validaci ve výuce. Doporučuje se zadávat ověřitelné zdroje a kontrolní otázky a učit studenty kritickému ověřování.
  \item[Bias / zkreslení] (general background knowledge) Sklony v datech nebo chování modelu, které mohou vést k nespravedlivým nebo nevhodným závěrům; při použití v kurzech je třeba pozornost při výběru a revizi obsahu. Zapojení různé perspektivy a revizní proces pomáhá odhalit a opravit zkreslení, a je vhodné mít transparentní auditní postupy.
  \item[GDPR a ochrana soukromí] (general background knowledge) Základní požadavky na zpracování osobních údajů (minimalizace, účelové omezení, anonymizace/pseudonymizace a dokumentace zpracování) — důležitá součást plánování nasazení AI ve výuce. Zahrnuje i informovaný souhlas studentů a bezpečné ukládání dat, včetně omezení přístupu a pravidelného vymazávání nepotřebných údajů.
\end{description}

Poznámka: položky označené jako "general background knowledge" jsou záměrně stručné; podrobné vysvětlení (včetně ilustrací tokenizace, vizualizace embeddings nebo RAG‑pipeline) najdete v odpovídajících kapitolách příručky. Tento slovník slouží jako startovací bod pro rychlou orientaci a plánování výuky a pro diskusi o praktických pravidlech nasazení v konkrétních třídních scénářích.